\section{Mapa de empatía y definición del problema}
\label{sec:mapa-empatia}

\subsection{Usuario y fuente}

El usuario representado es la persona que atiende la caja de un establecimiento de productos eléctricos en Quevedo. La información procede de un ensayo de entrevista estructurada realizado mediante juego de roles: Sánchez Centeno Roselyn Andreina interpretó a la cajera y Mendoza Palma Allan Jeremy formuló las preguntas. No se entrevistó a una persona que ejerciera efectivamente ese puesto; por ello, las respuestas constituyen material de exploración académica y no resultados de investigación de campo. El guion utilizado contiene trece preguntas (P1--P13), incluida la autorización de grabación (P1); el análisis temático se basa en P2--P13, documentadas en la sección~\ref{sec:ensayo-entrevista}.

El mapa sintetiza respuestas declaradas durante el ensayo e interpretaciones acotadas. Se utilizaron seis dimensiones: \emph{piensa y siente}, \emph{ve}, \emph{oye}, \emph{dice y hace}, \emph{esfuerzos} y \emph{resultados}. Cada elemento tiene una pregunta de origen, indicada en la figura~\ref{fig:mapa-empatia-cajera} y detallada en la tabla~\ref{tab:trazabilidad-empatia}. Esta organización es una síntesis del ensayo, no una medición psicológica de la participante.

\subsection{Mapa de empatía}

\begin{figure}[htbp]
\centering
\begin{tikzpicture}[
    x=\linewidth,
    y=1cm,
    celda/.style={draw=black!65,fill=black!3,rounded corners=3pt,
       text width=.425\linewidth,minimum height=3.15cm,
       inner sep=6pt,align=left,anchor=north,font=\footnotesize}
]
\node[celda] at (.25,0) {\textbf{PIENSA Y SIENTE (M01)}\par\smallskip
  Valora ahorrar tiempo en la atención (P13).\par\smallskip
  Prioriza evitar devoluciones por uso incorrecto del producto (P12).};
\node[celda] at (.75,0) {\textbf{VE (M02)}\par\smallskip
  Consulta códigos, referencias y fichas en el sistema (P4--P6, P11).\par\smallskip
  Describe clientes esperando cuando aumenta la demanda (P10).};
\node[celda] at (.25,-3.45) {\textbf{OYE (M03)}\par\smallskip
  Recibe solicitudes expresadas con nombres comunes de productos (P7, P13).\par\smallskip
  Se comunica por radio con bodega y consulta al personal técnico (P4, P11--P12).};
\node[celda] at (.75,-3.45) {\textbf{DICE Y HACE (M04)}\par\smallskip
  Registra productos, emite facturas y procesa pagos (P2, P5--P6).\par\smallskip
  Coordina existencias con bodega y solicita apoyo técnico (P3--P4, P12).};
\node[celda] at (.25,-6.90) {\textbf{ESFUERZOS (M05)}\par\smallskip
  Verifica transferencias fuera del sistema y busca productos de nombre técnico (P7).\par\smallskip
  Describe errores de cantidad o selección entre marcas y voltajes semejantes (P8--P9).};
\node[celda] at (.75,-6.90) {\textbf{RESULTADOS (M06)}\par\smallskip
  Desea búsquedas por sinónimos y nombres coloquiales (P13).\par\smallskip
  Propone validar transferencias con mayor eficiencia (P13) y evitar devoluciones por mal uso (P12).};
\end{tikzpicture}
\caption{Mapa de empatía del rol simulado de cajera, elaborado a partir del ensayo de entrevista. Los códigos M01--M06 remiten a la trazabilidad.}
\label{fig:mapa-empatia-cajera}
\end{figure}

La dimensión \emph{piensa y siente} registra prioridades inferidas de expresiones explícitas sobre tiempo y devoluciones, no emociones declaradas. \emph{Ve} y \emph{oye} describen información del entorno y comunicaciones mencionadas por la participante; no implican observación directa de un local ni transcripciones de conversaciones con terceros. El ensayo no permite determinar niveles reales de estrés, satisfacción, frecuencia de errores ni tiempos operativos.

\subsection{Hallazgos principales}

\noindent\textbf{H01. Desajuste entre el vocabulario del cliente y el catálogo.} En el ensayo se describió la búsqueda por nombre o código y la dificultad de encontrar productos cuyo nombre técnico no coincide con el nombre cotidiano utilizado por el cliente (P4, P7). La sugerencia de reconocer sinónimos (P13) refuerza la necesidad de identificar productos sin depender de una denominación exacta. Como interpretación, la diferencia de vocabulario puede prolongar la atención.

\noindent\textbf{H02. Riesgo de selección de productos con características semejantes.} La participante relató equivocaciones de cantidad y confusiones entre referencias de distinta marca o voltaje (P8), además de correcciones posteriores a la facturación (P9). También explicó que recurre a información técnica o a un vendedor especializado (P11--P12). Se identifica la necesidad de comprobar características relevantes antes del cobro, sin concluir que el cajero deba reemplazar al asesor técnico.

\noindent\textbf{H03. Interrupciones al confirmar transferencias.} La verificación manual mediante consulta fuera del sistema fue descrita como una actividad que consume tiempo (P7); la mejora de esa comprobación surgió espontáneamente al cierre (P13). Se infiere la necesidad de disponer de una confirmación oportuna y confiable antes de concluir una venta; el ensayo no permite fijar un tiempo máximo ni prescribir una integración bancaria.

\noindent\textbf{H04. Coordinación en momentos de alta demanda.} La participante señaló que, ante varios clientes esperando, pide a un vendedor apoyo anticipado para revisar precios y productos (P10), mientras que la consulta de existencias depende de bodega (P3--P4). Se interpreta que una coordinación clara puede contribuir a la continuidad de la atención. No se midieron filas, tiempos ni cargas de trabajo.

\subsection{Trazabilidad con la entrevista}

\begingroup
\small
\renewcommand{\arraystretch}{1.13}
\setlength{\LTpre}{6pt}
\setlength{\LTpost}{6pt}
\begin{longtable}{@{}>{\raggedright\arraybackslash}p{.09\linewidth}>{\raggedright\arraybackslash}p{.29\linewidth}>{\raggedright\arraybackslash}p{.42\linewidth}>{\raggedright\arraybackslash}p{.08\linewidth}@{}}
\caption{Origen de los elementos del mapa y de los hallazgos}\label{tab:trazabilidad-empatia}\\
\toprule
\textbf{ID} & \textbf{Dimensión o hallazgo} & \textbf{Sustento del ensayo} & \textbf{Origen} \\
\midrule
\endfirsthead
\toprule
\textbf{ID} & \textbf{Dimensión o hallazgo} & \textbf{Sustento del ensayo} & \textbf{Origen} \\
\midrule
\endhead
\bottomrule
\endlastfoot
M01 & Piensa y siente & Ahorro de tiempo considerado útil; propósito de evitar devoluciones por mal uso. & P12--P13 \\
M02 & Ve & Consulta de información del catálogo y descripción de clientes en espera. & P4--P6, P10--P11 \\
M03 & Oye & Solicitudes con nombres comunes; comunicación con bodega y personal técnico. & P4, P7, P11--P13 \\
M04 & Dice y hace & Registro, cobro, facturación, comprobación de existencias y consulta técnica. & P2--P6, P12 \\
M05 & Esfuerzos & Verificación manual, diferencia de nombres y errores con referencias parecidas. & P7--P9 \\
M06 & Resultados & Búsqueda por sinónimos, comprobación ágil de transferencias y prevención de mal uso. & P12--P13 \\
H01 & Identificación de productos & Búsqueda por código o nombre y discrepancia con la forma de pedir del cliente. & P4, P7, P13 \\
H02 & Exactitud técnica & Errores de cantidad o selección y dependencia de información especializada. & P8--P9, P11--P12 \\
H03 & Confirmación de pagos & Validación manual fuera del sistema y propuesta de mejorarla. & P7, P13 \\
H04 & Continuidad de atención & Coordinación con bodega y apoyo de vendedores durante la espera. & P3--P4, P10 \\
\end{longtable}
\endgroup

La trazabilidad diferencia lo \emph{declarado} en las respuestas de la \emph{interpretación} analítica. Las preguntas P2--P13 sustentan el mapa; P1 corresponde exclusivamente al consentimiento. No se identificaron observaciones directas del entorno de trabajo ni datos suficientes para cuantificar los problemas.

\subsection{Definición del problema}

Se consideraron tres formulaciones orientadas a las necesidades descritas:
\begin{enumerate}
\item \emph{Identificación:} ¿Cómo podríamos ayudar a quien atiende la caja a reconocer con precisión los productos solicitados, aunque los clientes utilicen nombres comunes o existan referencias técnicas similares?
\item \emph{Confirmación de pagos:} ¿Cómo podríamos facilitar que quien atiende la caja confirme oportunamente los pagos sin interrumpir la atención al cliente?
\item \emph{Alcance general:} ¿Cómo podríamos mejorar todos los procesos de venta de un establecimiento eléctrico?
\end{enumerate}

La primera opción se selecciona porque conecta las búsquedas por nombre (P4), las demoras por diferencias de vocabulario (P7), los errores entre referencias semejantes (P8) y la necesidad expresada al cierre (P13). La segunda también está respaldada (P7, P13), pero aborda una dificultad distinta que conviene estudiar de manera específica. La tercera es demasiado amplia para orientar un ejercicio de ideación a partir de este único ensayo.

\noindent\textbf{Definición seleccionada:}
\begin{quote}
\emph{¿Cómo podríamos ayudar a quien atiende la caja de un establecimiento eléctrico a identificar correctamente los productos solicitados, aun cuando el cliente utilice nombres coloquiales y existan referencias de características similares?}
\end{quote}

Esta pregunta se centra en la necesidad de la persona representada, puede investigarse mediante tareas observables de identificación y deja abiertas distintas alternativas de solución. No presupone inteligencia artificial, una interfaz específica ni una tecnología determinada. Dado que la fuente fue un juego de roles, la definición debe contrastarse con personal real del establecimiento antes de convertirla en requisitos del PFC \cite{pohl2025,ireb2026}.
